% !TEX root = ../main.tex

\begin{digest}
\noindent\textbf{Keywords:} code generation; execution feedback; reinforcement learning; supervised fine-tuning; PPO

\vspace{1em}

Code generation based on natural language instructions has become an important research topic in both artificial intelligence and software engineering. With the rapid development of large language models, modern code-oriented foundation models are able to read problem descriptions, infer algorithmic intentions, and generate executable programs in languages such as Python. These models are usually pretrained on massive code and text corpora, and then adapted to downstream tasks through supervised fine-tuning. This training paradigm has significantly improved the usability of code generation systems. However, code generation is different from ordinary text generation. A fluent and syntactically plausible answer is not necessarily a correct program. The final quality of generated code is determined by whether the program can be executed and whether it satisfies the functional requirements expressed by test cases. Therefore, there is a natural gap between token-level supervised learning objectives and execution-level correctness objectives.

This thesis studies reinforcement learning for code generation based on execution feedback. The central idea is to use the executability of code as a source of training signal. Instead of only asking the model to imitate a reference solution, the proposed workflow lets the model generate candidate programs, runs these programs against unit tests, converts the execution result into reward, and then uses reinforcement learning to optimize the generation policy. In this way, the training objective is more closely aligned with the evaluation criterion of code generation tasks. The work focuses on a small-scale but complete experimental setting: data processing, supervised fine-tuning, rollout collection, local execution, reward computation, advantage estimation, PPO policy update, and final evaluation are all implemented and connected into a full experimental loop.

The motivation of this work comes from two observations. First, a programming problem often has multiple correct solutions. Supervised fine-tuning on a single reference solution may encourage the model to reproduce a particular surface form, but it does not directly teach the model to satisfy the semantic constraints of the task. Second, execution feedback is relatively objective and easy to obtain for many code generation benchmarks. If a generated function passes all provided tests, it is more likely to be functionally correct; if it fails some tests, the number of passed tests and the error type can still provide useful partial feedback. This property makes code generation a suitable domain for reinforcement learning with automatic rewards.

The base model used in this thesis is Qwen2.5-Coder-1.5B, a code-oriented causal language model with strong pretrained programming ability. The experimental data are constructed from MBPP and HumanEval. MBPP is used for supervised fine-tuning, validation, and PPO rollout experiments, while HumanEval is used as the final cross-dataset test set. The raw datasets are converted into two unified formats. The SFT format contains problem prompts and reference responses, and is used for supervised fine-tuning. The RL format keeps the problem prompt, canonical solution, entry point, and test code, and is used for generation, execution, reward calculation, and policy optimization. After processing, the final dataset contains 385 MBPP training samples, 42 MBPP validation samples, and 164 HumanEval test samples.

Before training and evaluation, this thesis performs runtime validation on the processed datasets. For each sample, the canonical solution is executed together with the corresponding tests to ensure that the entry point, solution code, and test code are compatible with the local execution environment. This step is necessary because the reward signal in the PPO stage directly depends on test execution. If the test itself is malformed, or if the entry point cannot be resolved, the calculated reward would be unreliable. The runtime validation results show that all 591 processed samples can be correctly executed in the current framework. This provides a reliable data foundation for the following experiments.

The training pipeline starts with supervised fine-tuning. LoRA is used as the parameter-efficient fine-tuning method, so that the base model can be adapted under limited GPU memory. The SFT stage trains the model on 385 MBPP training samples for five epochs. The target modules include the major attention and feed-forward projection layers. The resulting SFT adapter is used as the baseline policy for subsequent PPO experiments. The SFT validation loss is 0.7472. Under the unified evaluation pipeline, the SFT baseline achieves a Pass@1 of 0.5714 on the MBPP validation set and 0.5244 on the HumanEval test set.

The core contribution of this thesis is the construction of an execution-feedback PPO loop for code generation. In the rollout stage, the current policy model receives a problem prompt and generates a candidate Python solution. Since language models may output explanations, Markdown code fences, repeated problem statements, or additional examples, the generated text must be cleaned before execution. The implemented post-processing module extracts the most plausible code block, aligns the function name with the required entry point when possible, and removes irrelevant text. The cleaned code is then executed in a local Python environment together with the test cases. The execution result includes whether the code is executable, how many tests are passed, whether all tests are passed, and whether an exception or assertion failure occurs.

Based on the execution result, the framework computes a scalar reward for each generated solution. The reward is then normalized into an advantage signal and used by PPO to update the LoRA adapter. The PPO implementation in this thesis is a simplified single-step version. Each generated program is treated as one episode, and the whole response is regarded as one action. The update objective includes the clipped policy loss, a scalar value loss, and a KL penalty against the reference policy. Although this is not a full token-level multi-step PPO implementation, it is sufficient for validating whether execution feedback can improve the generation behavior under limited data and computing resources.

Reward design is one of the key issues in this work. The thesis compares several reward functions. The original reward mainly uses the ratio of passed tests as the reward signal. This design is simple and directly related to functional correctness, but it may still be sparse when most generated solutions fail all tests. To provide denser feedback, \texttt{reward\_v2} combines three components: an executability reward, a partial test pass reward, and a full-pass bonus. This design distinguishes syntax or runtime failure, partial correctness, and complete correctness. It gives the model useful feedback even when the generated solution does not pass all tests.

This thesis also explores more complex reward designs, namely \texttt{reward\_v3} and \texttt{reward\_v3b}. These two variants introduce structural signals such as AST similarity and function signature matching. The intention is to check whether structural similarity to the reference solution can provide additional guidance beyond test execution. However, the experiments show that these structural rewards do not improve Pass@1 in the current setting. In some cases, they increase the average reward value but fail to improve the number of fully correct solutions. This indicates that, for small-scale code generation PPO, a more complex proxy reward may not be better if it is not well aligned with functional correctness.

The main validation results compare Base, SFT baseline, and the best PPO model on the MBPP validation set. The Base model achieves a Pass@1 of 0.6667, with 28 out of 42 samples fully passing the tests. The SFT baseline achieves a Pass@1 of 0.5714, with 24 fully passed samples. This result shows that small-scale supervised fine-tuning does not necessarily improve a strong code base model on every distribution. It may adapt the output style to the training data, but it can also damage part of the pretrained model's original capability. After PPO optimization with \texttt{reward\_v2}, the best PPO model achieves a Pass@1 of 0.6667, also with 28 fully passed samples. Therefore, PPO recovers the validation performance drop caused by SFT and matches the Base model on Pass@1.

The average reward results provide additional information. On the MBPP validation set, Base obtains an average reward of 0.7381, SFT obtains 0.6524, and PPO best obtains 0.7190. Although PPO best has the same Pass@1 as Base, its average reward is slightly lower. This suggests that Base may receive higher partial rewards on some failed samples, while PPO best more directly improves the number of fully passed samples. For code generation, Pass@1 is usually the more important metric because it measures whether a generated answer fully solves the problem in one attempt. Therefore, the validation result supports the conclusion that execution-feedback PPO can correct part of the SFT-induced performance degradation.

To test cross-dataset generalization, this thesis evaluates Base, SFT baseline, and PPO best on HumanEval. The Base model achieves a Pass@1 of 0.5000, the SFT baseline achieves 0.5244, and PPO best achieves 0.5122. The SFT model slightly improves over Base on HumanEval, while PPO best is between Base and SFT. The number of generation execution failures decreases from 7 for Base and 2 for SFT to 0 for PPO best, which shows that PPO improves the execution stability of generated outputs. However, PPO does not outperform SFT on HumanEval Pass@1. This result indicates that the PPO improvement mainly appears on the validation distribution used for rollout and tuning, and it does not fully transfer to the cross-dataset test setting.

The PPO hyperparameter experiments further reveal the sensitivity of reinforcement learning for code generation. The best original-reward PPO setting uses 42 rollouts, a learning rate of $1\times10^{-5}$, a KL coefficient of 0.02, and a batch size of 4. Reducing the KL coefficient to 0.01 leads to worse performance, because the policy moves too far away from the SFT policy and the mean probability ratio becomes too large. Reducing the learning rate to $5\times10^{-6}$ makes the update more conservative, but also weakens the improvement. Increasing the rollout number from 42 to 100 does not automatically improve performance either. These results suggest that PPO effectiveness depends on the balance among reward density, rollout scale, KL constraint, and update strength. More rollout data or a more aggressive update does not necessarily lead to better results.

The reward ablation experiments show that \texttt{reward\_v2} is the most effective design in this thesis. With \texttt{reward\_v2}, PPO best reaches a Pass@1 of 0.6667 on the MBPP validation set. In contrast, PPO with the original reward reaches 0.6190, while \texttt{reward\_v3} and \texttt{reward\_v3b} both reach 0.5952. The structural rewards in \texttt{reward\_v3} and \texttt{reward\_v3b} fail to improve Pass@1 even though they provide additional reward components. This finding is important because it shows that reward functions should not only be dense, but also aligned with the final evaluation target. For code generation, test execution remains the most reliable automatic signal in this experimental setting.

The case analysis provides a more concrete explanation of the model behavior. In a positive case from HumanEval, the task requires sorting the values at even indices of a list while keeping odd indices unchanged. The SFT output incorrectly uses an undefined variable, causing execution failure. The PPO output correctly extracts the even-indexed elements, sorts them, and writes them back to the corresponding positions. This case shows that execution feedback can help correct indexing and implementation logic errors. In another positive case, the task requires extracting fruit counts from a string. The SFT output uses a brittle hard-coded character access pattern, while the PPO output parses numeric tokens from the string and subtracts their sum from the total fruit count. This improves robustness for multi-digit quantities.

There are also negative cases. For example, in a HumanEval task that asks for the product of odd digits and requires returning 0 when all digits are even, Base and SFT correctly maintain a flag to handle the all-even case. PPO generates a shorter implementation based on reduction over odd digits. This implementation works when odd digits exist, but returns 1 instead of 0 when the input contains only even digits. This case shows that PPO can damage boundary condition handling. The reason is that the current PPO update is based on limited rollout data and simplified single-step credit assignment. If the reward does not sufficiently cover certain edge cases, the policy may learn a shortcut that passes many tests but fails important corner cases.

Overall, the experimental results support several conclusions. First, the complete SFT and PPO pipeline for execution-feedback code generation is feasible. The framework can process datasets, fine-tune a base model, generate code, execute tests, compute rewards, and update a policy adapter. Second, execution feedback can help repair the performance drop caused by small-scale SFT on the MBPP validation distribution. Third, reward design is critical. A direct and execution-aligned dense reward such as \texttt{reward\_v2} is more effective than more complicated structural rewards in the current setting. Fourth, PPO is sensitive to KL constraint, learning rate, rollout scale, and reward density. Fifth, the generalization improvement is still limited. PPO best improves validation Pass@1 relative to SFT, but does not surpass SFT on HumanEval.

This thesis also has several limitations. The most important limitation is data scale. The SFT training set contains only 385 MBPP samples, and the main PPO rollout set contains 42 validation samples in the best setting. Such a small scale is enough to verify the experimental pipeline, but it is insufficient for robust policy optimization. The second limitation is the simplified PPO formulation. Treating an entire generated program as one action greatly reduces implementation complexity, but it does not fully model token-level decisions and long-range credit assignment. The third limitation is reward coverage. The reward is based on available tests, so the model may overfit to the current test distribution if the tests are not comprehensive. The fourth limitation is the execution environment. The current local executor is suitable for MBPP and HumanEval style function-level tasks, but a stronger sandbox with timeout, resource isolation, and file system restrictions is required for more complex code generation scenarios.

Future work can improve this research in several directions. First, larger and more diverse training data should be introduced, including extended MBPP variants, EvalPlus, APPS, and real-world repository tasks. More diverse tasks may improve the cross-dataset generalization of PPO. Second, reward design can be improved by combining execution results with hidden tests, error classification, static analysis, boundary condition coverage, and code complexity constraints. Third, the PPO framework can be extended to token-level or multi-step trajectory optimization with a stronger value function and adaptive KL control. Fourth, the execution environment should be upgraded to a safer sandbox with timeout and resource control, so that generated code with infinite loops or unsafe operations can be handled robustly. Finally, evaluation should go beyond Pass@1 and average reward, and include Pass@k, code robustness, running time, readability, and failure type analysis.

In conclusion, this thesis verifies that execution feedback can be used as an effective reinforcement learning signal for code generation. Under small-scale data and limited computing resources, PPO based on execution feedback can improve the SFT model on the validation distribution and reduce execution failures. At the same time, the experiments reveal that strong code base models already have considerable capability, and that small-scale SFT and PPO do not automatically guarantee better generalization. The most practical lesson of this work is that execution-based reward is valuable, but it must be carefully designed and constrained. A simple, direct, and functionally aligned reward can be more effective than a complex proxy reward. This provides a concrete experimental reference for future research on reinforcement learning and execution-based optimization in code generation.
\end{digest}
